% !PRINTED 511
\weberpage{511}

Daraus folgt, daß $\varrho$ einer quadratischen Gleichung genügen muß, also daß
\weq{16}{\varrho=\pm i\quad\text{oder}\quad \varrho=\dfrac{\pm 1\pm\sqrt{-3}}{2}}
sein muß, weil dies die einzigen nicht reellen zwölften Einheitswurzeln sind, die einer quadratischen Gleichung genügen, und $\varrho=\pm1$ zu $x=x'$, $y=y'$ führen würde. Demnach muß $\sqrt D$ einem Körper angehören, der durch eine dieser Irrationalitäten $\varrho$ bestimmt ist, d.\,h. es muß
\weq{17}{D=-4m^2\quad\text{oder}\quad D=-3m^2}
sein. \gesperrt{Im ersten Fall} folgt aus (15)
\weq{18}{\begin{aligned}
y'&=\pm 2mx,\\
2mx'&=\mp y,
\end{aligned}}
also nach (8):
\[
n=m^2(x^2+x'^2),
\]
und da $n$ ungerade vorausgesetzt war, muß auch $m$ ungerade sein; nach (18) sind $y,y'$ und wegen (8) auch $x,x'$ gerade, und aus (12) folgt, da $\gamma$ und $\beta$ gerade sind:
% [sic] printed claim that $x,x'$ are both even conflicts with $n=m^2(x^2+x'^2)$ for odd $n,m$
\[
M^4=(\alpha+\beta\omega)^6=\left(\dfrac{y+ixm}{2}\right)^6,
\]
und aus (18):
\[
M'^4=(\alpha'+\beta'\omega)^6=\left(\dfrac{y'+ix'm}{2}\right)^6=-\left(\dfrac{y+ixm}{2}\right)^6,
\]
\weq{19}{M'^4=-M^4.}

Es haben also $M$ und $M'$ nicht gleiche, sondern entgegengesetzte 4te Potenzen.

Im zweiten der Ausnahmefälle (17) folgt in gleicher Weise aus (15):
\weq{20}{\begin{aligned}
\pm 2y'&=y-3mx,\\
\pm 2mx'&=y+mx.
\end{aligned}}

Darauf folgt, wenn man die oberen Zeichen nimmt,
\weq{21}{m(x'y-y'x)=2n,}
und hieraus schließt man, daß $m$ ungerade sein muß; denn wäre $m$ gerade, so müßte $y$ und $y'$ gerade sein, und nach (21) wäre $2n$ durch 4, also $n$ durch 2 teilbar. Dieser Fall ist also nicht weiter zu berücksichtigen, wenn wir ein gerades $D$ voraussetzen.

Hiermit ist folgender Satz bewiesen:

1. \gesperrt{Es sei $D\equiv 0\pmod 4$ eine negative Diskriminante, $x,y$ seien zwei Zahlen ohne gemeinsamen un-}
%CONT>>
